\documentclass[10pt,a4paper]{article}

\usepackage[a4paper,margin=2cm]{geometry}
\usepackage[utf8]{inputenc}
\usepackage[T1]{fontenc}
\usepackage{graphicx}
\usepackage{booktabs}
\usepackage{hyperref}
\usepackage{xurl}
\usepackage{amsmath}
\usepackage{enumitem}
\usepackage{xcolor}
\usepackage{titlesec}
\usepackage{parskip}
\usepackage{caption}

\hypersetup{
    colorlinks=true,
    linkcolor=blue,
    urlcolor=blue,
    citecolor=blue
}

\titleformat{\section}{\normalfont\large\bfseries}{\thesection}{1em}{}
\titleformat{\subsection}{\normalfont\bfseries}{\thesubsection}{1em}{}
\setlength{\parskip}{0.4em}

\title{\textbf{\large Information Retrieval Models — Vector Space and Probabilistic Approaches}}
\author{Muhammad Bagus Zulmi (25/572328/PPA/07191) \and Vicky Mahfudy (25/573019/PPA/07207)\\[0.3em]}
\date{}

\begin{document}

\maketitle

\vspace{-1.5em}

\section{Pendahuluan}

Laporan ini membandingkan performa 6 model \emph{ranked retrieval} pada korpus 14.343 dokumen berita berbahasa Indonesia (\texttt{News.csv}, kolom \texttt{content}): TF, TF-IDF, dan Word2Vec (\emph{TF-IDF weighted}) sebagai model \emph{vector space}; serta Query-Likelihood (QL) tanpa \emph{smoothing}, QL Laplace, dan QL \emph{linear interpolation} ($\lambda=0{,}5$) sebagai model probabilistik. Preprocessing yang diterapkan pada semua model: \emph{lowercasing}, penghapusan tanda baca, tokenisasi, dan penghapusan \emph{stopword} Bahasa Indonesia. Kelima query uji yang sama digunakan pada seluruh model, masing-masing mengambil \emph{top-10} dokumen. Kode dan hasil lengkap tersedia di \url{https://github.com/vickymahfudy/stbi-tugas-1}.

\section{Waktu Komputasi}

\begin{table}[h]
    \centering
    \small
    \begin{tabular}{lrrrr}
        \toprule
        \textbf{Model} & \textbf{Mean (s)} & \textbf{Std (s)} & \textbf{Min (s)} & \textbf{Max (s)} \\
        \midrule
        Word2Vec (TF-IDF weighted) & 0,0046 & 0,0023 & 0,0029 & 0,0086 \\
        QL (No Smoothing)          & 0,0087 & 0,0029 & 0,0060 & 0,0136 \\
        TF-IDF                     & 0,0130 & 0,0030 & 0,0101 & 0,0171 \\
        QL (Laplace)               & 0,0208 & 0,0057 & 0,0148 & 0,0301 \\
        QL (Linear Interpolation)  & 0,0248 & 0,0049 & 0,0163 & 0,0290 \\
        TF                         & 0,0624 & 0,0084 & 0,0569 & 0,0771 \\
        \bottomrule
    \end{tabular}
    \caption{Rata-rata waktu komputasi per model (5 query, korpus 14.343 dokumen).}
\end{table}

Word2Vec paling cepat karena \emph{similarity} dihitung pada ruang vektor berdimensi rendah (100) yang sudah dibangun sebelumnya. QL tanpa \emph{smoothing} juga cepat karena dokumen langsung tereliminasi (skor $-\infty$) begitu satu kata query tidak ditemukan. TF paling lambat karena implementasi \texttt{CountVectorizer} tidak melakukan \emph{precomputation} bobot seperti TF-IDF, membuat \emph{transform} dan \emph{cosine similarity} lebih mahal.

\section{Relevansi Dokumen Terhadap Query}

Sebagai proksi otomatis relevansi digunakan \texttt{term\_overlap\_ratio} (rasio kata query yang muncul literal di dokumen hasil), dirangkum pada Tabel~\ref{tab:relevansi}.

\begin{table}[h]
    \centering
    \small
    \begin{tabular}{lrr}
        \toprule
        \textbf{Model} & \textbf{Mean} & \textbf{Std} \\
        \midrule
        QL (No Smoothing)          & 1,000 & 0,000 \\
        QL (Linear Interpolation)  & 0,931 & 0,094 \\
        QL (Laplace)               & 0,909 & 0,113 \\
        TF                         & 0,781 & 0,190 \\
        TF-IDF                     & 0,754 & 0,203 \\
        Word2Vec (TF-IDF weighted) & 0,698 & 0,220 \\
        \bottomrule
    \end{tabular}
    \caption{Rata-rata proksi relevansi (\texttt{term\_overlap\_ratio}) per model.}
    \label{tab:relevansi}
\end{table}

Nilai 1,000 pada QL tanpa \emph{smoothing} bukan tanda relevansi terbaik, melainkan artefak mekanisme kerjanya: dokumen yang tidak memuat seluruh kata query otomatis dieliminasi, sehingga yang lolos filter pasti \emph{overlap} penuh, dan \emph{top-10}-nya seringkali tidak terisi penuh akibat \emph{data sparsity}. Sebaliknya, nilai Word2Vec yang paling rendah menandakan model ini menangkap kemiripan semantik, bukan berarti paling tidak relevan. Evaluasi relevansi yang valid membutuhkan penilaian manual terhadap isi dokumen, bukan hanya \emph{overlap} kata literal.

\section{Perbandingan Top-10 Antar Model}

\emph{Jaccard similarity} dari himpunan dokumen \emph{top-10} dihitung untuk setiap pasangan model pada masing-masing dari 5 query, lalu dirata-ratakan (Tabel~\ref{tab:jaccard}).

\begin{table}[h]
    \centering
    \small
    \begin{tabular}{lr}
        \toprule
        \textbf{Pasangan Model} & \textbf{Jaccard (rata-rata)} \\
        \midrule
        QL No Smoothing vs QL Linear Interp. & 0,579 \\
        TF vs QL Linear Interp.              & 0,271 \\
        TF vs TF-IDF                         & 0,251 \\
        TF-IDF vs QL Linear Interp.          & 0,247 \\
        TF vs Word2Vec                       & 0,209 \\
        QL No Smoothing vs QL Laplace        & 0,197 \\
        TF vs QL No Smoothing                & 0,172 \\
        QL Laplace vs QL Linear Interp.      & 0,168 \\
        TF-IDF vs Word2Vec                   & 0,161 \\
        TF-IDF vs QL No Smoothing            & 0,154 \\
        TF vs QL Laplace                     & 0,153 \\
        TF-IDF vs QL Laplace                 & 0,153 \\
        Word2Vec vs QL Linear Interp.         & 0,133 \\
        Word2Vec vs QL No Smoothing          & 0,098 \\
        Word2Vec vs QL Laplace               & 0,093 \\
        \bottomrule
    \end{tabular}
    \caption{Rata-rata \emph{Jaccard similarity} top-10 antar seluruh pasangan model (5 query).}
    \label{tab:jaccard}
\end{table}

QL tanpa \emph{smoothing} dan QL \emph{linear interpolation} paling mirip (Jaccard 0,579), wajar karena keduanya berbasis probabilitas kata yang sama sebelum \emph{smoothing} diterapkan. Word2Vec konsisten memiliki \emph{overlap} terendah dengan seluruh varian QL (0,09--0,13), mengonfirmasi representasi semantik menangkap relevansi dari sudut pandang berbeda dibanding model probabilistik berbasis kata literal.

\section{Efek Smoothing dan Normalisasi}

\subsection{Efek Smoothing pada Model QL}

Tabel~\ref{tab:smoothing} menunjukkan rata-rata jumlah dokumen yang berhasil dikembalikan (dari target \emph{top-10}) oleh masing-masing varian QL, dirata-ratakan atas 5 query.

\begin{table}[h]
    \centering
    \small
    \begin{tabular}{lrrr}
        \toprule
        \textbf{Varian QL} & \textbf{Rata-rata Hasil} & \textbf{Min} & \textbf{Max} \\
        \midrule
        Tanpa Smoothing      & 8,0  & 4  & 10 \\
        Laplace              & 10,0 & 10 & 10 \\
        Linear Interpolation & 10,0 & 10 & 10 \\
        \bottomrule
    \end{tabular}
    \caption{Rata-rata jumlah dokumen hasil (dari target top-10) per varian QL (5 query).}
    \label{tab:smoothing}
\end{table}

QL tanpa \emph{smoothing} pada 2 dari 5 query hanya mengembalikan sebagian dokumen (bahkan hingga 4 dari 10) karena mekanismenya mengeliminasi total dokumen yang tidak memuat \emph{seluruh} kata query (skor $-\infty$). \emph{Smoothing}, baik Laplace maupun \emph{linear interpolation}, menghilangkan masalah ini sepenuhnya: probabilitas kata yang tidak ada di dokumen tidak lagi nol, sehingga \emph{top-10} selalu terisi penuh pada seluruh query. Sebagai konsekuensinya, seperti terlihat pada Tabel~\ref{tab:relevansi}, proksi relevansi rata-rata sedikit menurun setelah \emph{smoothing} (1,000 $\rightarrow$ 0,909 untuk Laplace, 0,931 untuk \emph{linear interpolation}), karena dokumen yang dulunya tereliminasi (dan pasti \emph{overlap} penuh bila lolos) kini ikut masuk peringkat meski \emph{overlap} kata-nya tidak lengkap. \emph{Trade-off} ini wajar: \emph{smoothing} menukar sedikit presisi literal demi \emph{recall} yang jauh lebih baik.

\subsection{Efek Normalisasi (Stemming, Stopword Removal, Lowercasing)}

Seluruh model pada laporan ini sudah menerapkan \emph{lowercasing} dan penghapusan \emph{stopword} Bahasa Indonesia sebagai \emph{preprocessing} dasar (Bagian~1). Kami mencoba menambahkan langkah \emph{stemming} (Sastrawi) di atas model TF-IDF sebagai \emph{enhancement} lanjutan, namun proses \emph{stemming} pada \emph{vocabulary} unik dari 14.343 dokumen berjalan lebih dari 15 menit tanpa selesai (Sastrawi murni Python, lambat secara \emph{inherent} per kata) sehingga dihentikan sebelum menghasilkan data. Akibatnya, efek \emph{stemming} terhadap performa retrieval belum dapat dianalisis secara empiris pada laporan ini; evaluasi ini menjadi pekerjaan lanjutan, dengan kandidat solusi berupa \emph{caching} hasil \emph{stemming} ke \emph{disk} atau membatasi lingkup \emph{stemming} pada subset korpus yang lebih kecil.

\end{document}
